\documentclass[10pt,a4paper]{amsart}
\usepackage[margin=19mm]{geometry}
\usepackage{amsmath,amssymb,mathtools}
\usepackage[scr=rsfs,cal=euler]{mathalfa}
\usepackage{xfrac}
\usepackage[unicode,hidelinks,bookmarks=false]{hyperref}
\newcommand{\ie}{i.e.\ }
\newcommand{\cf}{cf.\ }
\newcommand{\resp}{respectively\ }
\newcommand{\Subsection}{\subsection}
\providecommand{\nobreakdash}{\nobreak\hbox{-}}
\setlength{\emergencystretch}{2em}
\multlinegap0pt
\usepackage{bm}
\usepackage{bbm}
\usepackage{dsfont}
\usepackage{xspace}
\usepackage{cancel}
\usepackage{mathtools}
\usepackage{amsthm}
\makeatletter
\@ifundefined{theo}{\newtheorem{theo}{Theo}}{}
\@ifundefined{lemma}{\newtheorem{lemma}[theo]{Lemma}}{}
\makeatother
\makeatletter
\expandafter\ifx\csname proof\endcsname\relax
\newenvironment{proof}{\noindent\textbf{Proof}.}{\hfill$\qed$}
\fi
\makeatother
\title[Lollipops, dense cycles and chords]{Lollipops, dense cycles and chords}
\author{Zden\v{e}k Dvo\v{r}\'ak, Beatriz Martins, St\'ephan Thomass\'e, Nicolas Trotignon}
\date{Source: arXiv:2502.04726v4 (CC BY 4.0)}
\begin{document}
\maketitle

\noindent\emph{Source and licence.} Lollipops, dense cycles and chords. Version arXiv:2502.04726v4 (CC BY 4.0), distributed under the Creative Commons Attribution 4.0 International licence, as recorded in the arXiv licence metadata for this exact version and shown on the abstract page https://arxiv.org/abs/2502.04726v4. Full rights notice: licenses/arxiv-2502.04726-CC-BY-4.0.txt.\par\smallskip

\noindent\textit{Editorial scope.} Counted unit: the Lemma whose source label is \texttt{l:contract} in the article (source0.tex lines 481--489, source label \texttt{l:contract}), delivered together with the article's own complete original proof of it (source0.tex lines 491--523, one \texttt{proof} environment of 33 lines). No mathematics is written, repaired or re-derived: statement and proof are the article's own text line for line. Required context retained uncounted: lem:activeneighborhood (lines 234--236); lem:kactive (lines 330--332); lem:nonactivepath (lines 470--472). Every definition, display and label of those lines is retained unchanged. Omitted from this package and not redistributed: 1--233, 237--329, 333--469, 473--480, 490--490, 524--980. Nothing inside a retained excerpt is omitted or altered: every retained body is delivered exactly as the article's lines are written, so no omission receipt of any kind accompanies this package. Citations: the retained excerpts carry no citation command at all, so no bibliography entry of the article is transcribed and no reference is redistributed. Reference adaptation: none was needed and none was made; every reference inside the delivered unit resolves inside it, so no unresolved reference or citation is produced. Printed slips kept literal: no printed word, symbol, hypothesis or number of the counted statement or of its proof was altered. Layout and class: the article's own class is \texttt{article} with packages including babel, amsmath, amssymb, amsthm, tikz, graphicx, hyperref, cleveref, amsfonts, titling, authblk; the wrapper is the lane's pinned \texttt{amsart} editorial wrapper at 10pt on A4, which restates the theorem-like environments the retained text needs and adds the article's own macro definitions that the retained text needs (none), with extra wrapper packages (bm, bbm, dsfont, xspace, cancel, mathtools). The delivered numbering is the unit's own. Assets: no retained span contains a figure environment, a graphics-inclusion command, a PDF-page inclusion, a table or a TikZ picture; no image file is copied into this package, the package declares no asset dependency and no image file is redistributed with it. Licence: version arXiv:2502.04726v4 of this article is distributed under the Creative Commons Attribution 4.0 International licence, as recorded in the arXiv licence metadata for this exact version and shown on the abstract page https://arxiv.org/abs/2502.04726v4. A full rights notice is delivered at licenses/arxiv-2502.04726-CC-BY-4.0.txt. Characters: every retained excerpt is pure ASCII, so the delivered engine, XeLaTeX with its default Latin Modern text font, reports no missing character.
\begin{lemma}\label{lem:activeneighborhood}
 If $G[C]$ contains a Hamiltonian path with ends $c_1$ and $u$, then $N_G(u) \subseteq V(C)$. 
\end{lemma}

\begin{lemma}\label{lem:kactive}
    $C$ contains at least $k$ active vertices. Moreover, if $d_C(c_1) < k$, then $C$ contains at least $k+1$ active vertices.   
\end{lemma}

\begin{lemma}\label{lem:nonactivepath}
    If $R$ is a subpath of $C$ containing only passive edges and $u$ is an active vertex, then $u$ has at most one neighbor in $R$.  Moreover, such a neighbor is an end of $R$. 
\end{lemma}

\begin{lemma}
    \label{l:contract}
    There exist sets $X_1,X_2\subseteq E(C)$ so that for $i\in\{1,2\}$, the graph $G_i$ obtained by deleting all vertices not in $C$ and contracting the edges of $X_i$
    has at least $k$ vertices and
    \begin{itemize}
        \item minimum degree at least $\left\lceil \frac{k+2}{2}\right\rceil$ if $i=1$, and 
        \item average degree at least $\tfrac{2}{3}(k+1)$ if $i=2$.
    \end{itemize}
\end{lemma}

\begin{proof}
    By~Lemma~\ref{lem:kactive}, $C$ has at least $k$ active vertices. Let $X_0$ be the set of 
    passive edges of $C$.  Let $G_0$ be the graph obtained from $G$ by deleting the vertices outside of $C$
    and contracting the edges of $X_0$; and let $C_0$ be the cycle in $G$ obtained from $C$ by this contraction.
    By~Lemma~\ref{lem:activeneighborhood} and Lemma~\ref{lem:nonactivepath}, the active vertices have the same degree in $G_0$ as in $G$.

    Note that the cycle $C_0$ is edgewise partitioned into edges whose both ends are active vertices, and
    paths of length~2 whose both ends are active and whose unique internal
    vertex is not.  Let $a_1b_1c_1$, \dots, $a_pb_pc_p$ be the paths
    of length~2 of $C_0$ such that the $a_i$'s and $c_i$'s are active and the
    $b_i$'s are not.  Suppose that $a_1, b_1, c_1$, \dots, $a_p, b_p, c_p$
    appear in this order along $C_0$ (note that possibly $c_i=a_{i+1}$,
    subscript taken modulo $p$).

    Let $X'_1=\{a_ib_i:i\in \{1,\ldots,p\}\}$ and $X_1=X_0\cup X'_1$; thus, $G_1$ is the graph obtained from $G_0$ by contracting the edges in $X'_1$.
    Note that the edges of $X'_1$ are vertex-disjoint.   Let $C_1$ be the cycle in $G_1$ obtained from $C_0$ by this contraction.
    Active vertices are in one-to-one correspondence to the vertices of $C_1$; in particular, $C_1$ has length at least $k$.
    Moreover, each vertex of $C_1$ has its two neighbors along $C_1$, and it is incident with at least $\left\lceil \frac{k-2}{2}\right\rceil$
    chords (because contracting the edges of $X'_1$ decreases the number of incident chords
    at most by half).  Therefore, $G_1$ has minimum degree at least $\left\lceil \frac{k+2}{2}\right\rceil$.

    We let $X_2=X_1$ if $G_1$ has average degree greater than $G_0$ and $X_2=X_0$ otherwise.  Thus $G_2 = G_1$ if $X_2=X_1$, and $G_2 = G_0$ otherwise.   Let $n_a$ be the number of edges in $E(G_0)\setminus E(C_0)$
    with both active ends, let $n_b$ be the number of edges in $E(G_0)\setminus E(C_0)$ with exactly one active end, and let $m$ be the number of active vertices.
    Since each active vertex is incident with at least $k-2$ chords in $G_0$,
    we have
    $$2n_a+n_b\ge (k-2)m.$$
    The average degree of $G_0$ is at least
    $$\frac{2|E(G_0)|}{|V(G_0)|}\ge \frac{2|C_0|+2n_a+2n_b}{|C_0|}=2+\frac{2n_a+2n_b}{|C_0|}\ge 2+\frac{n_a+n_b}{m}.$$
    The graph $G_1$ has at least $n_a$ chords and $m$ vertices, and thus it has average degree at least
    $$\frac{2|E(G_1)|}{|V(G_1)|}\ge 2+\frac{2n_a}{m}.$$
    Hence, $G_2$ has average degree at least
    $$2+\frac{\max(n_a+n_b,2n_a)}{m}\ge 2+\frac{\max((k-2)m-n_a,2n_a)}{m}\ge 2+\tfrac{2}{3}(k-2)=\tfrac{2}{3}(k+1).$$
\end{proof}

\end{document}
